\documentclass[journal]{./IEEE/IEEEtran}
\usepackage{cite,graphicx}

\newcommand{\SPTITLE}{LakBayanihan: A Community-Driven Multimodal Trip Planning Travel Companion and Real-Time Public Transportation Information System}
\newcommand{\ADVISEE}{Jonas R. Atienza}
\newcommand{\ADVISER}{Concepcion L. Khan}

\newcommand{\BSCS}{Bachelor of Science in Computer Science}
\newcommand{\ICS}{Institute of Computer Science}
\newcommand{\UPLB}{University of the Philippines Los Ba\~{n}os}
\newcommand{\REMARK}{\thanks{Presented to the Faculty of the \ICS, \UPLB\
                             in partial fulfillment of the requirements
                             for the Degree of \BSCS}}
        
\markboth{CMSC 190 Special Problem, \ICS}{}
\title{\SPTITLE}
\author{\ADVISEE~and~\ADVISER%
\REMARK
}
\pubid{\copyright~2024~ICS \UPLB}

%%%%%%%%%%%%%%%%%%%%%%%%%%%%%%%%%%%%%%%%%%%%%%%%%%%%%%%%%%%%%%%%%%%%%%%%%%

\begin{document}

% TITLE
\maketitle

% ABSTRACT
\begin{abstract}
Public transportation is a leading mode of transportation in the Philippines, but the public transportation system has often been criticized for lagging behind the standards of developed countries.
There is a lack of travel applications specifically tailored for public transportation assistance, and the current choices are limited and have limitations.
The goal of this paper is to develop the basis of a `commuter first' multimodal trip planning travel companion for mobile devices that commuters can use on demand and anywhere.
This is especially important for first-time commuters to an area.
The app will provide necessary real-time information about the current state of CPT modes on demand.
It will also help to empower and guide CPT operators and Local Government Units (LGUs) in giving accurate data.
Users will be able to suggest temporary and/or permanent changes, which will be subject to removal by authorized admins or through a reporting system.
This project has the potential to improve the Quality of Life (QoL) of many commuters, especially those who are first time commuters to a certain place.
It can also help commuters be aware of traffic congestions and make the proper decisions in exploring other possible and more efficient routes. 
\end{abstract}

% INDEX TERMS
\begin{keywords}
public transport, transportation assistance, gis, mobile development, react native, google maps, commuting, cross-platform, android, ios
\end{keywords}

% INTRODUCTION
\section{Introduction}
Transportation plays an important part in the life of the modern man. In the Philippines, public transportation is a leading mode of transportation---having a predominant number of registered Public Utility Vehicles (PUVs) as opposed to privately-owned Motor Vehicles (MVs) [INSERT UPDATED DATA/CITATION HERE;\ Current data from 2020, request from eFOI still pending]. Mass public transportation is also cheaper, relatively convenient, and can sometimes be faster than private modes of transportation~\cite{Narboneta16}.

However, the public transportation system in the Philippines has often been criticized for lagging behind the standards of developed countries, leading to a negative perception of its convenience and overall experience. Despite these challenges, the benefits of an efficient and well-organized public transport system cannot be ignored---both for the benefit of the people and for the benefit of a greener world~\cite{Lagua19}.

Currently, Metro Manila is the only city in the Philippines which has a formal public transport---namely, the Light Rail Transit (LRT) and the Metro Rail Transit (MRT). Other cities and municipalities are only mostly catered by buses and other privately-owned---and operated---informal modes of transportation such as \textit{jeepney}, van, \textit{`tricycle'} (auto rickshaws), \textit{`pedicab'} (cycle rickshaws), and \textit{`habal-habal'} (motorcycles-for-hire)~\cite{Mayo20}.
These---on some cities---are nowadays accompanied by the Modern Public Utility Vehicles (MPUVs) and Modern Public Utility Jeepneys (MPUJs) brought by the Public Utility Vehicle Modernization Program (PUVMP) of the Department of Transportation (DOTr) as issued in Department Order (D.O.) No. 2017--011. [ citation format for department order (?) ]

% NOTE: Check bib file for department order citation correctness.

\subsection{Background of the study}
Conventional Public Transport (CPT) is a mode of transport that is the backbone of transportation in the Philippines. It is recognized by law and by the public as the `standard' in the transport system. CPT covers the aforementioned \textit{jeepney}, van, \textit{`tricycle'}, \textit{`pedicab'}, \textit{`habal-habal'}, and train systems~\cite{RamizoJr19}.
`Conventional taxi'---such as the taxi systems in Metro Manila, Baguio City, etc.\---while under CPT, could be said to be distinct to the definition of CPT outlined in this paper, along with similar `ride hailing' services---such as GrabTaxi and Angkas, among others---since taxi services operate in a `pick up and drop off system' instead of adhering to a standard operating route.

The routes of each CPT mode have loops that are unique---albeit with some trips overlapping the routes of other trips. Each Trip T has a distinct average time to complete, and by extension, going from Point A to Point B on a different route might change a trip's duration---assuming the durations would be measured without the time it takes for drivers to wait for passengers.
This is of importance since time satisfaction, along with comfort and cost, plays a huge part in commuters' preference in choosing routes and modes for a commute~\cite{Mayo20}.

In the real world, unexpected events might trigger the unavailability of a CPT or its inability to operate at its full capacity. Events such as natural disasters, mobilizations, and transport strikes could be a reason for these inconveniences to the common commuter.
The location of various CPT transport terminals and their routes could also change over time. This might trigger confusion to commuters when asking locals that did not know of the said change since they would not be able to point them in the right direction. \\ \\ %\\ \\ \\ \\ 

% TODO: Find solution for rCol:footer overflow; [ \\ \\ \\ \\ \\ \\  from previous paragraph temporary ].

\subsection{Statement of the Problem}
Transport information (routes, cost, schedules, etc.) is not organized and available in a public database. While it is mandated by law for PUJs and other PUVs to display the \textit{`taripa'} (fare matrix) that they are adhering to, the only useful information for the commuters in it is the guide for the fares to the destinations under the route that they are in. 

Ride-hailing apps have been around in the Philippines for quite a while, as pioneered by GrabTaxi in 2013 and Uber in 2014. They were shortly followed by minor---turned major---players such as Angkas and Joyride~\cite{RamizoJr19}. While they are undoubtedly more convenient---and sometimes faster---than CPT, it is undeniable that the economical practicality of such is highly questionable~\cite{RamizoJr19}.
As such, it could be inferred that ride-hailing---or even the taxi services under CPT---is not a viable day-to-day option 

Currently, there is a lack of travel applications specifically tailored for public transportation assistance.
Since transportation through mass public transport is crucial in the lives of many Filipinos of differing backgrounds, an on demand and publicly accessible transport information database solution makes sense as a way to improve the Quality of Life (QoL) of many commuters, especially those who are first time commuters to a certain place.


\subsection{Objectives}
This paper aims to develop the foundations of a `commuter first' multimodal trip planning travel companion for mobile devices that commuters could access on demand and anywhere. It specifically aims to:

\begin{enumerate}
    \item Develop a client-facing mobile application for multimodal trip planning that would be a travel companion and as such, provide necessary real-time information about the current state of CPT modes on demand.
    \item Develop a management application for web and mobile that would be able to make necessary real-time updates to the information system.
    \item Empower and guide CPT operators and the Local Government Units (LGUs) in giving accurate data.
    \item Let regular users suggest temporary and/or permanent changes that would be subjected to removal---if deemed necessary---by either authorized admins (CPT operators and LGU personnel) or through a reporting system.
    \item Assess and evaluate the usefulness, ease of use, and satisfaction of the users of the system.
    \item Lay the proper foundation for further open source development of the project as a possible national effort.
\end{enumerate}

\subsection{Significance of the Study}
Prefering mass public transport in general is the best choice in going from one point to another [ Insert Sakay.ph about page citation ].
It's a novel solution to the country's long-standing traffic congestion problem that is the worst on---but is not unique to---Metro Manila~\cite{RamizoJr19}.
Consulting the current rankings of the iOS App Store and Google Play Store as of November 2023, current choices would be limited to a small pool including Google Maps, Waze, and Sakay.ph. However, both Google Maps and Waze are focused on giving driving directions and while Google Maps has support for bus and train trip schedules/terminals, it still has its own limitations.
Sakay.ph, on the other hand, is a promising solution. However, its usability is limited to the National Capital Region (NCR). Observing their tagline `Your Metro Manila Commuting Guide', it can be inferred that regional support is of low priority. The website's FAQs page states that there is already early data added for Bacolod, Baguio, Batangas City, Cebu City, Davao, General Santos City, Iloilo, Laguna, and Legaspi City.
It is also stated that only jeepney, bus, and walk itineraries are included in the said data [ Insert SAKAY.PH website citation ].
It also falls short in User Interface (UI) design and User Experience (UX) when subjected to evaluation using UX and UI principles [ Insert UI/UX citation; see ref 9, 10 in FSAM paper (guide) ].
Licensing the project using a copyleft open source license may provide other developers with a feature-rich API, as well as an opportunity to expand the
project to harbor the needs of the whole country for a publicly available transport information database. 

% REVIEW OF RELATED LITERATURE
\section{Review of Related Literature}
A number of works related to this study has been implemented in various transportation networks. Almost every implementation has been catered to the benefit of private transportation modes or formal CPTs such as rail systems and buses.
Although each implementation is unique in terms of what framework and methodology they use, their end goal is the same---which is to provide commuters and travelers [on private vehicles] valuable information regarding the trips that they are about to make.

Gkiotsalitis and Stathopoulosi~\cite{Gkiotsalitis15} wrote `A Mobile Application for Real-Time Multimodal Routing Under a Set of Users’ Preferences' which introduced a mobile navigation application designed for users without private vehicles or those preferring a combination of private and public transport to minimize travel time, offering intermodal routes tailored to individual preferences within complex urban networks, supported by a real-time database and two algorithms for optimal path determination.
A study made by Noreikis and his team~\cite{Noreikis14} created an optimal in-vehicle park-and-ride solution for seamless private-public transport mode transfer which used multimodal trip planning technologies.
Dimokas and his colleagues~\cite{Dimokas18} introduced the combination of already widely available carpooling solutions and public transport modes in their application.

The Philippines, however, has a scarcity of these types of applications. While it is true that Google Maps---being the leading interactive maps application in the country---has a public transport route planning feature under Google Transit, it is very limited.
Google Transit only supports LRT and MRT routes, as well as some bus networks in NCR\@.
As of writing this paper, there is no support yet for \textit{jeepneys} and other CPT modes---or even signs of Google supporting it in the future.
Sakay.ph is a notable application produced during the Philippine Transit App Challenge in 2013 which is successful in implementing a localized version of Google Transit which also accomodates \textit{jeepney} routes.
The application has long been updated to accomodate P2P, UV Express, and Pasig River Ferry routes. The usability of Sakay.ph is limited to NCR and some major cities such as Cebu and Davao, according to its website. However, some users have reported that some route suggestions are impractical and is not commonly used by them as noted in a paper by Narboneta and Teknomo~\cite{Narboneta15}.
A study conducted by Comandao and her colleagues~\cite{Comandao21} in their paper titled `E-Jeep NN\@: An Android E-Jeep Guide App with GPS and Open Trip Planner API for the New Normal' had promising results in terms of usability and acceptance of their application.
However, it seems that their application, `E-Jeep NN' is not available for public use.

% MATERIALS AND METHODS
\section{Materials and Methods}
This chapter covers the tools utilized for application development, the incorporation of features, and the assessment of the application's implementation.

\subsection{Development Environment}
\subsubsection{\textbf{Host Machine}}
A laptop with the following specifications was used to host the development environment of the application:

\begin{itemize}
    \item $\textbf{Operating System:}$  Arch Linux x86\_64
    \item $\textbf{Host:}$ ASUS TUF Gaming A15 FA506IV\_FX506IV 1.0
    \item $\textbf{CPU:}$ AMD Ryzen 9 4900H @ 3.300GHz
    \item $\textbf{Memory:}$ 31511MiB $\approx$ 32GB DDR4
\end{itemize} \hfill \

\subsubsection{\textbf{Client Machines}}
Since the aim of the study is to produce a mobile application for both Android and iOS hosts, two (2) devices were used for testing and debugging of the application:
\begin{description}
    \item[\textbf{Android Client:}] \hfill \\
        A Xiaomi Poco F3 with the following specifications was used:
        \begin{itemize}
            \item $\textbf{Operating System:}$ Android 14 aarch64
            \item $\textbf{Host:}$ Xiaomi Poco F3 5G M2012K11AG
            \item $\textbf{Chip:}$ Qualcomm Snapdragon 870 @ 3.190GHz 
            \item $\textbf{Memory:}$ 5622MB $\approx$ 6GB
            \item $\textbf{Display:}$ 1080px$\times$2400px @ 392ppi 
        \end{itemize}
    \item[\textbf{iOS Client:}] \hfill \\
        An iPhone XS with the following specifications was used:
        \begin{itemize}
            \item $\textbf{Operating System:}$ iOS 17.1.1
            \item $\textbf{Host:}$ Apple iPhone XS A2097
            \item $\textbf{Chip:}$ Apple A12 Bionic @ 3.190GHz 
            \item $\textbf{Memory:}$ 3775MB $\approx$ 4GB
            \item $\textbf{Display:}$ 1125px$\times$2436px @ 458ppi 
            \end{itemize}
\end{description}

\hfill

\subsubsection{\textbf{Development Tools}}
The following tools were used in developing the application:

\begin{description}
    \item[\textbf{Vim:}] \hfill \\
        A terminal-based text editor that served as the main development environment of the application.

    %\item[\textbf{Visual Studio Code:}] \hfill \\
    %    A graphical text editor that served as an alternative development environment for some tasks that Vim cannot do.

    \item[\textbf{Git:}] \hfill \\
        A version control system that is used to monitor progress and each change in the development process of the application.

    \item[\textbf{React:}] \hfill \\
        A JavaScript library for frontend web UI development that was used to develop the management application on desktop. 

    \item[\textbf{React Native:}] \hfill \\
        A JavaScript library for cross-platform UI development in mobile devices that was used to develop the client-facing mobile application.

    \item[\textbf{Expo:}] \hfill \\
        A set of tools built on top of React Native that is used as an intermediary [server] between the development host machine and the development client machines.

    \item[\textbf{Expo Go:}] \hfill \\
        A mobile client for Expo that runs development builds.

    \item[\textbf{General Transit Feed Specification (GTFS):}] \hfill \\
        An open standard for distributing information about transit systems that served as the schema of the transit system database.
    
    \item[\textbf{Leaflet:}] \hfill \\
        An open-source JavaScript library for mobile-friendly interactive maps which served as the interactive map of the application.
    
    \item[\textbf{complexity-report:}] \hfill \\
        An npm package for software complexity analysis of JavaScript projects which benchmarked the complexity of the source code of the application.
\end{description}

\subsection{Feature set}
The application is composed of a client-facing application and a management application. \hfill \\

\subsubsection{\textbf{Client-facing application}}

\begin{description}
    \item[\textbf{Register:}] \hfill \\
        Clients are required to register in order to use extended functionalities[*] of the application.

    \item[\textbf{Generate multiple routes:}] \hfill \\
        Multiple routes will be generated from going from point A to point B with predefined user preferences (preferred CPT modes) being shown first.

    \item[\textbf{Generate public alerts*:}] \hfill \\
        Registered clients can generate public alerts (using markers or cirlces) such as vehicular accidents to alert commuters. These alerts are moderated in an upvote/downvote system and can be taken down by administrators.

    \item[\textbf{Add Routes*:}] \hfill \\
        Registered clients can add routes or request changes in a particular route. A prompt asking for GTFS-ready information will be used.
\end{description}

\subsubsection{\textbf{Management application}}

\begin{description}
    \item[\textbf{Login:}] \hfill \\
        Administrators are required to login using a provided account.

    \item[\textbf{View routes:}] \hfill \\
        Administrators can visualize the GTFS data being used by the application.

    \item[\textbf{Add and edit routes:}] \hfill \\
        Administrators can make changes to the GTFS database.

    \item[\textbf{Approve/Reject requests:}] \hfill \\
       Administrators can approve or reject requests made by clients.
    
    \item[\textbf{Moderate alerts:}] \hfill \\
       Administrators can create, modify, and take down public alerts.
\end{description}


\subsection{User Evaluation}
For the usability and user experience aspect as denoted in objectives 3, 4, and 5, the application is tested by students residing in Batangas Province, transport operators, and transport enthusiasts willing to take part in the study. For evaluators using an Android device, the APK of the application will be provided. For iPhone users, Expo Go will be used to communicate with the Expo server that will be ran by the proponent.

The usability and user experience aspect will be measured using the Usefulness, Satisfaction, and Ease of Use (USE) Questionnaire~\cite{Lund01}, a 7-point likert scale, which will be modified to accommodate objectives 3 and 4. 

For the maintainability aspect as denoted in objective 6, the application is tested by senior BS Computer Science students, the UPLB ICS Faculty, and other IT experts. The GitHub repository used in the development process of the application will be provided for them. A software complexity report generated by the npm package `complexity-report' will also be provided.

In measuring the maintainability of the source code, a questionnaire will be used to measure the estimated rebuild value, percentage of redundant code, lines of code per unit, cyclomatic complexity per unit (provided by the report), number of parameters per unit, and number of incoming calls per module.
The metrics would be used to evaluate the volume, duplication, unit size, unit complexity, unit interfacing, and module coupling properties that would be mapped to the ISO/IEC 9126 metric.
Additional information in the form of remarks will also be in the questionnaire which will be subjected to qualitative assessment.

% RESULTS AND DISCUSSION
\section{Results and Discussion}
The quick brown fox jumps over the lazy dog. The quick brown fox jumps over
the lazy dog. The quick brown fox jumps over the lazy dog. The quick brown
fox jumps over the lazy dog.

% CONCLUSION AND FUTURE WORK
\section{Conclusion and Future Work}
The quick brown fox jumps over the lazy dog. The quick brown fox jumps over
the lazy dog. The quick brown fox jumps over the lazy dog. The quick brown
fox jumps over the lazy dog.

% APPENDICES
\appendices

\section{Proof of the First Zonklar Equation}
Appendix one text goes here...

\section{}
Appendix two (without title) text goes here...

% ACKNOWLEDGMENT
\section*{Acknowledgment}
Many thanks to...

% BIBLIOGRAPHY
\bibliographystyle{./IEEE/IEEEtran}
\bibliography{./atienza-sp-ieee}
% \nocite{*}

% BIOGRAPHY
\begin{biography}[{\includegraphics{./yourPicture.eps}}]{Jonas R. Atienza}
Text here\ldots
\end{biography}


\end{document}
 
